% =============================================================================
% Kapitel 5: Auswertung
% =============================================================================
\chapter{Ergebnisse}
\label{ch:auswertung}

% ─────────────────────────────────────────────────────────────────────────────
\section{Parameterstudie}
\label{sec:ergebnisse-paramscan}

Die systematische Variation der Fit-Parameter über 83~Kombinationen
ergab klare Trends, die in Abbildung~\ref{fig:param-scan} zusammengefasst
sind.

\begin{figure}[htbp]
  \centering
  \TODO{Hier fig8\_parameter\_scan.png und fig9\_parameter\_heatmap.png
        einfügen.}
  %\includegraphics[width=\textwidth]{fig8_parameter_scan.png}
  %\includegraphics[width=\textwidth]{fig9_parameter_heatmap.png}
  \caption{Ergebnisse der Parameterstudie.
    Oben: Residuen aller 83~Läufe, sortiert nach Gesamtresidual.
    Unten: Heatmaps der besten Residuen für Kombinationen aus
    $k$-Gewichtung und Vergleichsraum (links), $k_{\min}$ und
    $N_{\text{legs}}$ (Mitte) sowie kantenaufgelöste Residuen der
    zehn besten Läufe (rechts).}
  \label{fig:param-scan}
\end{figure}

Die wesentlichen Befunde sind:
\begin{itemize}[nosep]
  \item Der \textbf{Wavelet-Raum} liefert durchgehend niedrigere
        Residuen als der $k$- oder $R$-Raum, da er simultane Empfindlichkeit
        in beiden Dimensionen bietet.
  \item \textbf{$k^2$-Gewichtung} stellt den besten Kompromiss dar;
        $k^1$ unterschätzt die Beiträge schwerer Rückstreuer,
        $k^3$ verstärkt das Hochfrequenzrauschen.
  \item \textbf{Mehrfachstreuung} ($N_{\text{legs}} = 4$) verbessert
        den Fit gegenüber reiner Einfachstreuung ($N_{\text{legs}} = 2$)
        signifikant.
  \item Der Einfluss von $k_{\min}$ (\num{2} vs.\ \num{3}) ist gering;
        $k_{\min} = 3$ wird bevorzugt, da der Bereich $k < 3$ stärker
        durch Hintergrundartefakte kontaminiert ist.
  \item In der zweiten Phase zeigt sich, dass größere Schrittweiten
        (\SI{0.005}{\Angstrom}) und ein Clusterradius von \SI{6}{\Angstrom}
        die besten Ergebnisse liefern.
\end{itemize}

% ─────────────────────────────────────────────────────────────────────────────
\section{Probe~1: CoNiCu-Mittelentropielegierung}
\label{sec:probe1}

\subsection{Messbedingungen}

\begin{table}[htbp]
  \centering
  \caption{Messbedingungen und Datenbereich der MEA-Probe.}
  \label{tab:probe1-bedingungen}
  \small
  \begin{tabular}{lll}
    \toprule
    & \textbf{Wert} & \textbf{Anmerkung} \\
    \midrule
    Probe          & \TODO{CoNiCu MEA} & \TODO{Probenbeschreibung} \\
    Kanten         & Co-K, Ni-K, Cu-K & Simultane Messung \\
    Energiebereich & \SIrange{7500}{10000}{\electronvolt} & \\
    $k$-Bereich    & \SIrange{3}{11.5}{\per\Angstrom} & Alle drei Kanten \\
    \TODO{Temperatur} & \TODO{RT / \SI{300}{\kelvin}} & \\
    \TODO{Messgeometrie} & \TODO{Transmission / Fluoreszenz} & \\
    \bottomrule
  \end{tabular}
\end{table}

\subsection{Strukturvergleich: FCC vs.\ BCC vs.\ HCP}

Tabelle~\ref{tab:struct-comparison} vergleicht die Fitresiduen der
drei Ausgangsstrukturen.

\begin{table}[htbp]
  \centering
  \caption{Kantenaufgelöste und gewichtete Gesamtresiduen für
    verschiedene Ausgangsstrukturen der MEA-Probe.
    Der niedrigste Gesamtresidual ist hervorgehoben.}
  \label{tab:struct-comparison}
  \small
  \begin{tabular}{lccccccc}
    \toprule
    \textbf{Struktur} &
    $R_{\text{Co}}$ & $R_{\text{Ni}}$ & $R_{\text{Cu}}$ &
    $R_{\text{total}}$ &
    $S_0^2$(Co) & $S_0^2$(Ni) & $S_0^2$(Cu) \\
    \midrule
    FCC & \TODO{0.000} & \TODO{0.000} & \TODO{0.000} & \TODO{\textbf{0.000}} &
          \TODO{0.00} & \TODO{0.00} & \TODO{0.00} \\
    BCC & \TODO{0.000} & \TODO{0.000} & \TODO{0.000} & \TODO{0.000} &
          \TODO{0.00} & \TODO{0.00} & \TODO{0.00} \\
    HCP & \TODO{0.000} & \TODO{0.000} & \TODO{0.000} & \TODO{0.000} &
          \TODO{0.00} & \TODO{0.00} & \TODO{0.00} \\
    \bottomrule
  \end{tabular}
\end{table}

\TODO{Interpretation einfügen: Welche Struktur liefert den besten Fit?
Ist das Ergebnis konsistent mit der Literatur (FCC erwartet für CoNiCu)?
Wie groß ist der Unterschied?}

Abbildung~\ref{fig:struct-comparison} zeigt den direkten Vergleich der
EXAFS-Spektren und Fouriertransformierten aller drei Strukturen.

\begin{figure}[htbp]
  \centering
  \TODO{Hier fig10\_structure\_comparison.png einfügen.}
  %\includegraphics[width=\textwidth]{fig10_structure_comparison.png}
  \caption{Strukturvergleich der MEA-Probe.
    (a)~Residuen nach Kante und Struktur.
    (b)--(d)~$\chi(k)\cdot k^2$ für Co, Ni, Cu mit experimentellen
    Daten (schwarz) und Fits aus FCC (grün), BCC (blau), HCP (rot).
    (e)--(g)~Zugehörige Fouriertransformierte.
    (h)~Tabellarische Zusammenfassung der Fitparameter.}
  \label{fig:struct-comparison}
\end{figure}

\subsection{Konvergenz}

\begin{figure}[htbp]
  \centering
  \TODO{Hier Konvergenzplots einfügen (struct\_fcc\_convergence.png etc.).}
  %\includegraphics[width=0.7\textwidth]{struct_fcc_convergence.png}
  \caption{Konvergenzverlauf des Gesamtresiduals für den FCC-Fit der
    MEA-Probe.}
  \label{fig:convergence-mea}
\end{figure}

\TODO{Beschreibung des Konvergenzverhaltens: Ab welcher Iteration
stagniert der Residual? Wie vergleichen sich die drei Strukturen?}

\subsection{EXAFS-Fit und Fouriertransformation}

\begin{figure}[htbp]
  \centering
  \TODO{Hier struct\_fcc\_exafs\_ft.png einfügen.}
  %\includegraphics[width=\textwidth]{struct_fcc_exafs_ft.png}
  \caption{EXAFS-Fit der MEA-Probe (FCC-Struktur).
    Obere Reihe: $\chi(k)\cdot k^2$ (berechnet vs.\ experimentell)
    für Co, Ni und Cu.
    Untere Reihe: Zugehörige Fouriertransformierte.
    Der grau hinterlegte Bereich ($k < \SI{3}{\per\Angstrom}$) wurde
    vom Fit ausgeschlossen.}
  \label{fig:exafs-ft-mea}
\end{figure}

\TODO{Qualitative Beschreibung: Wo stimmt der Fit gut überein? Gibt es
systematische Abweichungen? Welche Kante zeigt den höchsten Residual
und warum (z.\,B.\ geringeres Signal-zu-Rausch-Verhältnis bei Cu)?}

\subsection{Radiale Verteilungsfunktionen}

\begin{figure}[htbp]
  \centering
  \TODO{Hier struct\_fcc\_rdf.png einfügen.}
  %\includegraphics[width=\textwidth]{struct_fcc_rdf.png}
  \caption{Partielle radiale Verteilungsfunktionen $g_{ij}(R)$ der
    MEA-Probe für die drei Absorbertypen.
    Jedes Panel zeigt die Beiträge der drei Nachbarelemente.}
  \label{fig:rdf-mea}
\end{figure}

\TODO{Beschreibung der RDF-Peaks: Position des ersten Peaks
(Nächste-Nachbar-Abstand), Breite (Unordnung), relative Höhe
(Koordinationszahl).}

\subsection{Koordinationszahlen und Bindungslängen}

\begin{table}[htbp]
  \centering
  \caption{Aus der \evax{}-Analyse bestimmte Koordinationszahlen und
    mittlere Bindungslängen der MEA-Probe (erste Koordinationsschale).
    \TODO{Werte nach Abschluss der Fits einfügen.}}
  \label{tab:cn-mea}
  \small
  \begin{tabular}{lSSS}
    \toprule
    \textbf{Paar} &
    {$N_{ij}$} &
    {$\langle R_{ij}\rangle$ (\si{\Angstrom})} &
    {$\sigma_{ij}$ (\si{\Angstrom})} \\
    \midrule
    Co--Co & \TODO{0.0} & \TODO{0.00} & \TODO{0.00} \\
    Co--Ni & \TODO{0.0} & \TODO{0.00} & \TODO{0.00} \\
    Co--Cu & \TODO{0.0} & \TODO{0.00} & \TODO{0.00} \\
    \midrule
    Ni--Co & \TODO{0.0} & \TODO{0.00} & \TODO{0.00} \\
    Ni--Ni & \TODO{0.0} & \TODO{0.00} & \TODO{0.00} \\
    Ni--Cu & \TODO{0.0} & \TODO{0.00} & \TODO{0.00} \\
    \midrule
    Cu--Co & \TODO{0.0} & \TODO{0.00} & \TODO{0.00} \\
    Cu--Ni & \TODO{0.0} & \TODO{0.00} & \TODO{0.00} \\
    Cu--Cu & \TODO{0.0} & \TODO{0.00} & \TODO{0.00} \\
    \bottomrule
  \end{tabular}
\end{table}

\TODO{Interpretation: Stimmen die Gesamtkoordinationszahlen
($\sum_j N_{ij} \approx 12$ für FCC)? Sind die Bindungslängen
konsistent mit dem Vegard'schen Gesetz?}

\subsection{Chemische Nahordnung}

\begin{table}[htbp]
  \centering
  \caption{Warren-Cowley-Parameter $\alpha_{ij}$ der ersten
    Koordinationsschale der MEA-Probe.
    Negative Werte zeigen Paarbildungstendenz an, positive Werte
    Vermeidung.}
  \label{tab:sro-mea}
  \small
  \begin{tabular}{lccc}
    \toprule
    & Co & Ni & Cu \\
    \midrule
    Co & \TODO{$+0.00$} & \TODO{$-0.00$} & \TODO{$\phantom{+}0.00$} \\
    Ni & \TODO{$-0.00$} & \TODO{$+0.00$} & \TODO{$\phantom{+}0.00$} \\
    Cu & \TODO{$\phantom{+}0.00$} & \TODO{$\phantom{+}0.00$} & \TODO{$+0.00$} \\
    \bottomrule
  \end{tabular}
\end{table}

\TODO{Interpretation: Welche Paare werden bevorzugt/vermieden? Ist die
Legierung ein idealer Mischkristall ($\alpha \approx 0$) oder gibt es
lokale Ordnung? Stimmen die Ergebnisse mit DFT-Vorhersagen oder
anderen experimentellen Befunden überein?}

% ─────────────────────────────────────────────────────────────────────────────
\section{Probe~2: CoPt-Nanopartikel}
\label{sec:probe2}

\subsection{Messbedingungen}

\begin{table}[htbp]
  \centering
  \caption{Messbedingungen und Datenbereich der CoPt-Probe.}
  \label{tab:probe2-bedingungen}
  \small
  \begin{tabular}{lll}
    \toprule
    & \textbf{Wert} & \textbf{Anmerkung} \\
    \midrule
    Probe          & \TODO{CoPt-Nanopartikel} & \TODO{Beschreibung} \\
    Kanten         & \TODO{Co-K, Pt-L$_3$} & \\
    Energiebereich & \TODO{...} & \\
    $k$-Bereich    & \TODO{...} & \\
    \TODO{Temperatur} & \TODO{...} & \\
    \bottomrule
  \end{tabular}
\end{table}

\subsection{Strukturvergleich}

\TODO{Analog zur MEA-Probe werden hier FCC, BCC und ggf.\ L1$_0$
verglichen. Für CoPt-Nanopartikel ist die Frage besonders relevant,
ob eine geordnete Phase (L1$_0$) vorliegt oder eine statistische
Verteilung auf dem FCC-Gitter.}

\begin{table}[htbp]
  \centering
  \caption{Strukturvergleich der CoPt-Probe.}
  \label{tab:struct-comparison-copt}
  \small
  \begin{tabular}{lcccc}
    \toprule
    \textbf{Struktur} &
    $R_{\text{Co}}$ & $R_{\text{Pt}}$ & $R_{\text{total}}$ &
    $S_0^2$(Co) \\
    \midrule
    FCC  & \TODO{0.000} & \TODO{0.000} & \TODO{0.000} & \TODO{0.00} \\
    L1$_0$ & \TODO{0.000} & \TODO{0.000} & \TODO{0.000} & \TODO{0.00} \\
    \TODO{...} & & & & \\
    \bottomrule
  \end{tabular}
\end{table}

\subsection{EXAFS-Fit und Fouriertransformation}

\begin{figure}[htbp]
  \centering
  \TODO{EXAFS- und FT-Plots der CoPt-Probe einfügen.}
  \caption{EXAFS-Fit der CoPt-Probe.
    \TODO{Analoges Layout wie Abb.~\ref{fig:exafs-ft-mea}.}}
  \label{fig:exafs-ft-copt}
\end{figure}

\subsection{Koordinationszahlen und Bindungslängen}

\begin{table}[htbp]
  \centering
  \caption{Koordinationszahlen und Bindungslängen der CoPt-Probe.}
  \label{tab:cn-copt}
  \small
  \begin{tabular}{lSSS}
    \toprule
    \textbf{Paar} &
    {$N_{ij}$} &
    {$\langle R_{ij}\rangle$ (\si{\Angstrom})} &
    {$\sigma_{ij}$ (\si{\Angstrom})} \\
    \midrule
    Co--Co & \TODO{0.0} & \TODO{0.00} & \TODO{0.00} \\
    Co--Pt & \TODO{0.0} & \TODO{0.00} & \TODO{0.00} \\
    Pt--Co & \TODO{0.0} & \TODO{0.00} & \TODO{0.00} \\
    Pt--Pt & \TODO{0.0} & \TODO{0.00} & \TODO{0.00} \\
    \bottomrule
  \end{tabular}
\end{table}

\subsection{Chemische Nahordnung}

\begin{table}[htbp]
  \centering
  \caption{Warren-Cowley-Parameter der CoPt-Probe.}
  \label{tab:sro-copt}
  \small
  \begin{tabular}{lcc}
    \toprule
    & Co & Pt \\
    \midrule
    Co & \TODO{$+0.00$} & \TODO{$-0.00$} \\
    Pt & \TODO{$-0.00$} & \TODO{$+0.00$} \\
    \bottomrule
  \end{tabular}
\end{table}

\TODO{Interpretation: Deutet der SRO-Parameter auf L1$_0$-Ordnung hin
($\alpha_{\text{Co-Pt}} < 0$, $\alpha_{\text{Co-Co}} > 0$)?
Oder handelt es sich um eine statistische Legierung?}

% ─────────────────────────────────────────────────────────────────────────────
\section{Vergleich beider Proben}
\label{sec:vergleich}

\TODO{Dieser Abschnitt soll die Ergebnisse beider Proben
gegenüberstellen. Mögliche Vergleichsaspekte:}

\begin{itemize}[nosep]
  \item Grad der chemischen Nahordnung ($\alpha$-Parameter) in beiden
        Systemen,
  \item Fitqualität in Abhängigkeit von der Anzahl der Komponenten
        (ternär vs.\ binär),
  \item Lokale Verzerrung (Streuung der Bindungslängen) als Maß für
        den Gitterverspannungseffekt,
  \item Robustheit der Ergebnisse gegenüber der Wahl der
        Ausgangsstruktur.
\end{itemize}

\begin{table}[htbp]
  \centering
  \caption{Zusammenfassender Vergleich der beiden Probensysteme.}
  \label{tab:vergleich}
  \small
  \begin{tabularx}{\textwidth}{lXX}
    \toprule
    & \textbf{CoNiCu MEA} & \textbf{CoPt NP} \\
    \midrule
    Kristallstruktur  & \TODO{FCC}  & \TODO{FCC / L1$_0$} \\
    Gesamtresidual    & \TODO{0.000} & \TODO{0.000} \\
    Nächster-Nachbar-Abstand & \TODO{\SI{0.00}{\Angstrom}} &
                               \TODO{\SI{0.00}{\Angstrom}} \\
    Gesamt-KZ (1.~Schale) & \TODO{$\approx 12$} & \TODO{$\approx 12$} \\
    Stärkster SRO-Effekt  & \TODO{Co-Ni bevorzugt} &
                             \TODO{Co-Pt bevorzugt} \\
    \bottomrule
  \end{tabularx}
\end{table}
